\section{Síntesis y ampliación}

\subsection{Guía de escenarios aplicables para cada método}

\begin{table}[H]
\centering
\begin{tabular}{p{0.22\textwidth}p{0.35\textwidth}p{0.33\textwidth}}
\toprule
\textbf{Escenario} & \textbf{Método recomendado} & \textbf{Justificación} \\
\midrule
Reparación de imágenes & Manipulación del Score (RePaint) o DPS & La operación con máscara es sencilla; ambas clases de métodos son viables \\
Superresolución & Guiado por pseudoinversa & El submuestreo es una operación lineal, por lo que se puede utilizar la solución analítica \\
Desenfoque gaussiano & Guiado por pseudoinversa & La convolución es una operación lineal \\
Reconstrucción de CT/MRI & Guiado por pseudoinversa & La operación de proyección es lineal \\
Recuperación de distorsión no lineal & DPS & Se requiere manejar el operador no lineal $\mathcal{A}$ \\
Generación con control de layout & Manipulación del Attention Map & Es necesario controlar las relaciones espaciales correspondientes \\
Transferencia de estilo & Manipulación del Score & Implica el intercambio de características entre imágenes \\
\bottomrule
\end{tabular}
\caption{Guía para la selección de métodos de guiado durante la inferencia en diferentes escenarios de aplicación}
\end{table}

\subsection{Repaso del contenido central de esta clase}

Esta clase presenta sistemáticamente el marco y los métodos básicos de \textbf{generación guiada durante la inferencia}, abarcando los siguientes contenidos centrales:

\begin{enumerate}
    \item \textbf{Motivación}: Aprovechar las características de generación iterativa de los modelos de difusión preentrenados para inyectar señales de guía durante la inferencia, logrando así una generación condicional sin necesidad de volver a entrenar.
    \item \textbf{Manipulación del Score}: Métodos prácticos que manipulan directamente el score de des ruido o los mapas de atención (como RePaint), sin garantías teóricas pero ampliamente efectivos.
    \item \textbf{Marco de problemas inversos}: Modelar la generación condicional como un problema inverso $y = \mathcal{A}(x) + n$, proporcionando herramientas de análisis sistemáticas.
    \item \textbf{Descomposición bayesiana}: $\nabla_{x_t} \log p(x_t|y) = \nabla_{x_t} \log p(x_t) + \nabla_{x_t} \log p(y|x_t)$, reduciendo el problema a la aproximación del score de verosimilitud.
    \item \textbf{Guiado por pseudoinversa}: Basado en una aproximación gaussiana, aplicable a mapeos hacia adelante lineales con una precisión de aproximación relativamente buena.
    \item \textbf{DPS}: Basado en una aproximación Delta, aplicable a cualquier mapeo hacia adelante diferenciable; su implementación es extremadamente simple pero la aproximación es más rugosa.
\end{enumerate}

\subsection{Conexión con la siguiente clase}

Esta clase cubre las dos primeras de tres categorías principales de tecnologías de guiado durante la inferencia. La próxima clase (Lecture 12: Inference-Time Guided Generation 2) discutirá tecnologías más avanzadas de la tercera categoría, que podrían incluir:
\begin{itemize}
    \item Métodos de aproximación del score de verosimilitud más precisos
    \item Guiado durante la inferencia basado en optimización (como DDRM, MCG, etc.)
    \item Aplicaciones específicas del guiado durante la inferencia en dominios como la generación 3D y de video
\end{itemize}

\subsection{Lecturas adicionales}

\begin{itemize}
    \item \textbf{RePaint}: Lugmayr et al., ``RePaint: Inpainting using Denoising Diffusion Probabilistic Models,'' CVPR 2022.
    \item \textbf{DPS}: Chung et al., ``Diffusion Posterior Sampling for General Noisy Inverse Problems,'' ICLR 2023.
    \item \textbf{DDRM}: Kawar et al., ``Denoising Diffusion Restoration Models,'' NeurIPS 2022.
    \item \textbf{Classifier Guidance}: Dhariwal \& Nichol, ``Diffusion Models Beat GANs on Image Synthesis,'' NeurIPS 2021.
    \item \textbf{Sobre problemas inversos y modelos de difusión}: El orador recomienda consultar el paper de revisión (survey) para conocer más sobre técnicas de manipulación del score y resolución de problemas inversos.
\end{itemize}

%% --- Fin del contenido --- %%

\end{document}
